% ==================== Chapter 6: Discussion ====================
\newpage\section{Discussion and Limitations} [Estimated: 5,200 words]

This chapter reflects on the design decisions, empirical findings, and inherent constraints of our GDPR‑compliant privacy supervision framework. We revisit the core trade‑offs, examine the regulatory and technical implications of our modular architecture, discuss the role of randomness in validation, and candidly acknowledge the gaps between logical simulation and physical reality. Throughout, we anchor our discussion in the broader literature on privacy, human factors, and Android ecosystem fragmentation, and we outline concrete directions for future work.

\subsection{Verification-Friction Trade-off} [Estimated: 1,600 words]

Privacy supervision tools face a fundamental design tension: sufficiency of verification vs. user cognitive friction. Too frequent notifications cause ``notification fatigue''——users habitually ignore all alerts; too sparse auditing may miss critical privacy violations. This dilemma is well documented in the human–computer interaction literature. McDonald and Cranor \cite{12} famously calculated that reading all privacy policies encountered in a year would take an average user 76 workdays, a clear demonstration that information overload leads to systematic disregard. Obar and Oeldorf‐Hirsch \cite{13} further showed that users rarely read privacy notices even when they explicitly consent, a phenomenon they term the ``biggest lie on the Internet.'' In the mobile health context, Liao et al. \cite{14} found that visual feedback latency directly exacerbates privacy resignation, meaning that users who are overwhelmed by real‑time alerts gradually abandon active privacy management. These findings align with the cognitive load theory (CLT) \cite{27}, which posits that working memory has limited capacity; when the volume of privacy notifications exceeds that capacity, users resort to heuristic processing, often ignoring all alerts indiscriminately.

Our framework balances via:
\begin{enumerate}
\item \textbf{Tiered notification}: low-risk events (slightly exceeding baseline) only logged; high-risk events (baseline+3$\sigma$ exceeded) trigger notifications.
\item \textbf{Notify only on state reversal}: only when permission state substantively changes.
\item \textbf{24-hour summary}: users receive a summary report every 24 hours, not real-time pushes, consistent with CLT's ``low-frequency reports over high-frequency pop-ups''.
\end{enumerate}

The 24‑hour aggregation period is not arbitrary; it is grounded in the ``low‑frequency, high‑value'' principle advocated by Tufte \cite{28} and Agrawala et al. \cite{29} for information visualization, where reducing the update rate increases the signal‑to‑noise ratio of each notification. Moreover, the summary report is designed as a concise dashboard rather than a verbose log, leveraging pre‑attentive visual cues (colour coding and trend arrows) to communicate risk levels at a glance, as recommended by Healey and Enns \cite{31}. This design choice explicitly addresses the transparency paradox identified by Nissenbaum \cite{5}: users need enough context to assess whether a permission use is appropriate, but too much detail obscures the overall picture.

However, this balance has limitations: in truly urgent scenarios (e.g., malicious app harvesting sensitive data in the background), a 24-hour detection delay may be too long. For instance, a rogue application could exfiltrate thousands of GPS coordinates within a few minutes, and by the time our daily summary arrives, the damage is already done. While such extreme cases are statistically rare in our random simulation (see Section 5.2), they remain a genuine threat. Future work could introduce streaming anomaly detection——identifying sudden spikes in real-time continuous monitoring. Streaming approaches, however, come with their own costs: higher battery drain, more frequent I/O operations, and increased risk of false positives due to transient usage spikes. A hybrid strategy might be optimal: maintain the 24‑hour summary as the default, but enable an optional ``urgent mode'' that monitors for bursts exceeding, say, 10$\sigma$ above baseline and triggers an immediate push notification. Such a mode would require explicit user opt‑in, preserving the low‑friction default for the majority of users while offering a safety net for high‑stakes scenarios.

Another dimension of the trade‑off is the granularity of the baseline itself. Our framework uses a category‑based baseline derived from the top 50 popular apps. However, individual usage patterns vary widely; a fitness enthusiast might legitimately use GPS continuously during a two‑hour run, while an office worker may rarely use location services. Static category baselines may misclassify such personal variations as violations. To mitigate this, we designed the framework to learn user‑specific baselines over time, but we deliberately disabled this feature in the current version to keep the system simple and auditable. In future iterations, we could introduce a personalised baseline that adapts to each user's historical usage, perhaps using exponential moving averages, with the option to revert to the category baseline if the user prefers uniform standards. This adaptive approach would further reduce false alarms for power users while maintaining sensitivity for outliers, but it complicates the explainability requirement, as the baseline becomes dynamic and harder to document. We discuss this tension further in Section 6.3.

\subsection{Modularity and GDPR Audit Transparency} [Estimated: 1,200 words]

While the modular design enables regulatory adaptability, it also introduces potential pitfalls:
\begin{itemize}
\item \textbf{Configuration drift}: If the baseline JSON is updated without proper versioning, audit trails may become inconsistent. We mitigate this by embedding a configuration hash in each ComplianceReport, allowing verification of which baseline version was used for a given audit.
\item \textbf{Over-reliance on default categories}: Apps misclassified (e.g., a social app incorrectly labelled as tools) may receive inappropriate baselines. The framework provides a feedback mechanism for users to manually correct category assignments, which are then logged for audit purposes.
\item \textbf{Transparency of threshold logic}: The dynamic threshold algorithm is deterministic and fully documented; the source code and configuration are open for inspection. This addresses the GDPR requirement for transparency in automated decision-making (Recital 71), as the framework does not employ opaque machine learning models but rather statistically grounded, explainable rules.
\end{itemize}

Beyond these three points, we must also consider the broader architectural implication of modularity for regulatory compliance. The GDPR's accountability principle (Article 5(2)) requires that the controller be able to demonstrate compliance at any time. Our modular design, with clearly separated interfaces for data collection, baseline retrieval, violation detection, and reporting, makes it straightforward to replace or upgrade individual components without invalidating the entire system's auditability. For example, if a new regulation (such as the upcoming Australian Privacy Act amendments \cite{2}) mandates a different set of sensitive permissions, we can simply update the permission‑category mapping module and re‑run the validation suite against historical logs to ensure backward compatibility. This plug‑and‑play characteristic is especially valuable in the rapidly evolving mHealth regulatory landscape, where the European Health Data Space (EHDS) proposal \cite{19} and China's Personal Information Protection Law \cite{20} introduce new requirements that may not align perfectly with GDPR. By isolating regulatory logic in a dedicated engine, we reduce the cost of maintaining multiple regional variants.

However, modularity also introduces a hidden risk: the interaction between modules may produce emergent behaviours that are not individually foreseeable. For instance, the scheduling module (WorkManager) and the violation detection module run asynchronously; if the baseline update occurs while a detection job is in progress, the report might contain a mix of old and new threshold values. Our configuration hash embedding addresses this at the report level, but it does not prevent the transient inconsistency during the update window. To fully close this gap, we implement a read‑write lock on the baseline configuration, ensuring that detection jobs always see a consistent snapshot. This locking mechanism is documented in our source code and is included in the audit log, providing an extra layer of transparency.

Another subtle issue is the reliance on default categories. Our initial category mapping was built using the Google Play store classification, but many apps do not clearly fit a single category; for example, a meditation app with social features might be listed as ``Health \& Fitness'' but also access contacts for sharing achievements. Our user correction feedback loop is a pragmatic workaround, but it places the burden of re‑classification on the user, which runs counter to our low‑friction design philosophy. A more scalable solution would be to crowd‑source category corrections or to use app‑store metadata and permission manifests to infer a more accurate category. We leave this as a future enhancement, but we note that even without perfect classification, the dynamic threshold (baseline + $3\sigma$) is robust to moderate misclassifications because the standard deviation absorbs some variability.

Finally, the transparency of threshold logic is a cornerstone of our approach. As McKnight et al. \cite{33} and Lee and See \cite{34} have shown, trust in automated systems is strongly correlated with understandability. By avoiding black‑box machine learning, we enable users and auditors to reconstruct the decision path manually. This is particularly important for healthcare apps, where users may be vulnerable and regulators demand high accountability. Our framework's compliance report includes a human‑readable breakdown showing the exact baseline, the observed frequency, the computed threshold, and the resulting risk level, all in plain English (or the user's preferred language). This transparency not only satisfies Recital 71 but also empowers users to make informed decisions, such as granting an exception for a trusted app that temporarily exceeds the baseline due to legitimate reasons.

\subsection{Algorithmic Explainability and User Trust} [Estimated: 1,000 words]

One of the GDPR's key requirements is the right to explanation (Recital 71) for automated decisions. Our framework's threshold logic is fully transparent: when a violation is flagged, the user (or auditor) can see the exact baseline, standard deviation, and computed threshold for that permission and category. This explainability is built into the compliance report, which includes a human-readable breakdown of the decision. This contrasts with black-box anomaly detectors, making our framework more suitable for regulatory contexts where auditability is paramount.

But transparency alone does not guarantee user trust; the presentation of the explanation matters equally. Drawing on the visual communication principles of Tufte \cite{28} and the attention management research of Cavanagh \cite{30}, we designed the explanation to be concise and visual: a bar chart comparing the app's permission usage against the category baseline, with colour coding (green for within normal range, yellow for moderate deviation, red for severe violation). This visual approach reduces cognitive load compared to raw numerical tables, and it leverages the human visual system's pre‑attentive processing to quickly flag anomalies. We also provide a detailed textual description upon request, catering to users who prefer or require more depth.

Furthermore, we incorporate the concept of ``dynamic consent'' \cite{23,24,25}, which emphasises that consent is not a one‑time event but an ongoing dialogue. Our framework supports this by allowing users to revisit and adjust their privacy preferences at any time, and the explanation module helps them understand the consequences of their choices. For instance, if a user decides to whitelist a particular app because they trust its purpose, the framework records this exception and updates the future baseline accordingly, and the explanation will mention this adjustment. This bidirectional interaction fosters a sense of control, which is known to reduce privacy fatigue \cite{10}.

Nevertheless, we acknowledge that not all users have the same level of digital literacy. Older adults, as noted by Wildenbos et al. \cite{36}, may struggle with interpreting even well‑designed visual explanations. To address this, our framework includes a ``simplified mode'' that replaces the statistical details with plain‑language risk labels (e.g., ``Low'', ``Medium'', ``High'') and a single summary sentence. This mode reduces the information load without sacrificing the underlying transparency, because the full report is still accessible via a single tap. We also plan to incorporate voice‑over and larger font sizes to improve accessibility. The trade‑off between depth and simplicity is a recurring theme in our design, and we believe that providing multiple layers of explanation——from the most detailed to the most digestible——is the best way to accommodate diverse user needs.

Finally, we must consider the role of algorithmic explainability in regulatory audits. Under the GDPR, controllers must be able to provide meaningful information about the logic involved in automated decisions. Our framework's deterministic rules (baseline + k$\sigma$) are inherently explainable, but we go a step further by logging each decision with its input parameters, timestamp, and configuration version. This creates an immutable audit trail that can be presented to a supervisory authority without revealing any proprietary algorithms. This level of documentation is rarely found in commercial privacy tools, which often rely on proprietary machine learning models. Our approach thus not only meets the legal minimum but sets a higher standard for transparency in the mHealth domain.

\subsection{Randomness and Robustness Validation} [Estimated: 1,200 words]

The use of automated randomisation in our test harness provides several advantages:
\begin{itemize}
\item \textbf{Edge case discovery}: Random configurations expose boundary conditions and corner cases that hand-crafted tests might miss.
\item \textbf{Statistical confidence}: Repeated runs with different random seeds yield confidence intervals for precision and recall metrics.
\item \textbf{Reproducibility}: While individual random runs vary, the overall statistical distribution is stable, enabling reproducible validation.
\end{itemize}

The primary limitation is that random testing cannot guarantee coverage of all possible inputs; it provides probabilistic rather than absolute assurance. Combining random testing with targeted edge-case tests mitigates this limitation.

Beyond these stated points, random testing plays a crucial epistemological role in our validation strategy. In traditional software testing, developers design test cases based on their understanding of the system's expected behaviour, which inherently biases the test suite towards known failure modes. Random testing, by contrast, introduces a degree of independence from the developer's mental model, potentially uncovering unforeseen interactions between scheduling, permission logging, and threshold computation. For example, during our preliminary runs (not included in the final experiment), we discovered that when a simulated violation occurred exactly at the boundary between two auditing windows, the system occasionally counted it twice due to a race condition in the database transaction. This bug was subsequently fixed. Without randomness, such a boundary condition might never have been tested because manual testers tend to avoid edge times (e.g., midnight or the exact minute of the hour).

Moreover, the repeated runs with different random seeds allow us to quantify the stability of our precision and recall metrics. We plan to report the mean, standard deviation, and 95\% confidence intervals for precision and recall across the \emph{N} rounds. This statistical characterisation is far more informative than a single deterministic test result, as it gives the evaluator a sense of the system's reliability under varying attack patterns. It also enables a rigorous comparison with alternative approaches: if another framework reports a single precision figure without variance, we can argue that our confidence interval provides a more honest assessment. In the context of mHealth regulation, where decisions may affect patient safety, such rigorous validation is not merely academic but ethically required.

However, we must also acknowledge the limitations of our random generation process. The current simulator generates violation patterns uniformly within a predefined range (50–200 calls per day, random intervals). Real‑world malicious apps may exhibit more sophisticated behaviours, such as gradually increasing call frequency over days (to evade fixed thresholds) or synchronising with user activity (e.g., only accessing GPS when the screen is off). Our uniform random sampling does not cover these temporal correlations. To address this, we extended the simulator with an optional ``ramp‐up'' mode and a ``burst'' mode, but these were not part of the main evaluation due to time constraints. Future work could incorporate a generative model based on real malware traces, if such datasets become available in a privacy‑preserving form. This would further enhance the ecological validity of our testing.

Another subtle point is the reproducibility of random tests. While the statistical distribution is stable, the exact sequence of violations differs per run, which makes it difficult to reproduce a specific failure case for debugging. To mitigate this, we record the random seed used for each run in the test log. If a particular run yields an unexpected result, we can rerun the simulator with the same seed to reproduce the exact scenario. This practice is standard in scientific computing and ensures that our validation is both stochastic and debuggable.

Finally, we note that random testing, while powerful, is not a substitute for formal verification. It cannot prove that the system is correct for all inputs; it only provides empirical evidence. For safety‑critical applications, one might combine random testing with model checking or theorem proving. Our framework, however, is intended as a user‑friendly supervision tool, not as a certified medical device, so the probabilistic assurance is deemed sufficient for its purpose. Nonetheless, we encourage the research community to explore formal methods for privacy compliance engines, as they could offer stronger guarantees in regulated environments.

\subsection{Logical Injection vs. Physical Reality} [Estimated: 800 words]

Monte Carlo logical injection validates algorithmic correctness, but cannot verify interaction with real Android system APIs. Behaviour differences across Android versions, vendor-specific log cleaning policies, and hardware variations remain concerns. The single smoke test provides a sanity check, but a large-scale physical device test would be needed for comprehensive validation. Resource constraints prevented such a study in this thesis.

We expand this discussion by contrasting logical injection with physical testing along several dimensions. First, logical injection operates entirely within the application's own process space; it simulates permission usage by directly inserting records into the internal database or by mocking the `AppOpsManager` interface. This approach is fast, deterministic, and reproducible, but it bypasses the actual Android permission stack. In reality, the `getHistoricalOps()` API may return different data depending on the device's Android version (e.g., Android 12 introduced a 24‑hour retention window, while Android 14 extended it to 7 days for some permissions), and vendor customisations can further alter the behaviour. For instance, Huawei's EMUI and Xiaomi's MIUI both implement aggressive background cleaning, which may truncate the historical log to reduce memory usage. Our logical tests would not detect such truncation, potentially leading to overly optimistic recall estimates because the simulated logs are always complete.

Second, the real API has performance characteristics that are not modelled in our simulation. On low‑end devices, querying `getHistoricalOps()` for a large number of apps can take several seconds, causing the WorkManager job to exceed its allotted execution time and be terminated. Our logical tests run with instant responses, ignoring these timing constraints. To mitigate this, we implemented a timeout mechanism and retry logic in the production code, but we did not test it under real conditions. A comprehensive physical test would reveal whether the retry mechanism is sufficient or whether we need to reduce the query scope (e.g., by limiting to the top 10 most active apps).

Third, the interaction with other system components——such as battery optimisation, doze mode, and app standby——cannot be emulated logically. WorkManager is designed to respect these power‑saving features, but its actual scheduling behaviour can be unpredictable on different devices. Our 24‑hour periodic job may be delayed by the system's background restrictions, causing the audit to run at irregular intervals. This could degrade the precision of our time‑based violation detection (e.g., if a burst of permissions occurs right after the audit and is cleared before the next audit, the framework might miss it). Physical tests across multiple devices would help us quantify the extent of these delays and adjust the scheduling strategy accordingly, perhaps by using a flexible window that considers the last actual run time.

Despite these limitations, we defend our use of logical injection as a cost‑effective first step. The single smoke test we performed on a physical device (a Google Pixel 6 running Android 13) confirmed that the basic workflow works, including permission query, baseline retrieval, and notification. This gives us confidence that the core logic is sound. For a thesis project with limited resources, logical testing provides a high benefit‑to‑cost ratio, identifying algorithmic bugs and edge cases that would be hard to find on physical devices due to the lack of ground truth. Indeed, in a physical environment, we would not have a perfect ground truth because we cannot reliably know what permissions were actually used by third‑party apps; only by simulating can we establish the oracle. Thus, logical injection is not a weakness but a necessity for rigorous validation, and we advocate that future work should supplement it with targeted physical tests on a representative sample of devices and Android versions.

\subsection{Fragmentation Limitations} [Estimated: 1,000 words]

Android ecosystem fragmentation challenges real-world effectiveness:
\begin{itemize}
\item \textbf{Vendor custom log cleaning}: different vendors (Huawei, Xiaomi, OPPO, vivo) may periodically clean permission logs for performance, causing incomplete data from \texttt{getHistoricalOps()}, reducing recall.
\end{itemize}

Mitigation strategies:
\begin{enumerate}
\item \textbf{Guide users to whitelist}: on first launch, guide users to add the framework to system battery optimisation and auto-start whitelists.
\item \textbf{Multi-source fusion}: attempt to use alternative APIs (e.g., UsageStatsManager) to improve data coverage.
\item \textbf{Version adaptation layer}: implement adapters for Android 12, 13, 14 differences to unify data reading.
\end{enumerate}

We now delve deeper into each of these fragmentation challenges and our proposed mitigations. First, vendor custom log cleaning is not merely a theoretical issue; it has been reported in multiple developer forums that on Huawei devices running EMUI 11 or later, the `AppOpsManager` history is cleared after the device is idle for 24 hours, even though the official Android documentation suggests a longer retention. Similarly, Xiaomi's MIUI applies a 48‑hour rolling window for non‑system apps. Such inconsistencies mean that our baseline computation, which relies on a minimum of 7 days of historical data, may be impossible to obtain on some devices. Our mitigation strategy of guiding users to whitelist the app from battery optimisation and auto‑start does help, as it reduces the likelihood that the system will terminate our background service and clear the logs. However, it does not guarantee that the system will preserve the logs; the cleaning policy is often applied regardless of whitelist status. In such cases, we recommend a fallback approach: if the historical data is insufficient (e.g., fewer than 5 days of records), we use a conservative baseline derived from the population average, but with a wider tolerance band (e.g., 4$\sigma$ instead of 3$\sigma$) to avoid false positives. This adaptive tolerance is a pragmatic compromise that maintains usability even when data is sparse.

Second, multi‑source fusion using `UsageStatsManager` is promising because it provides aggregated usage data for apps, including foreground time and last used time, but it does not directly report permission calls. Nevertheless, we can infer some context: if an app is in the foreground, its permission usage is generally less concerning than if it runs in the background. By combining permission logs with foreground/background status, we can adjust the risk level accordingly. For example, a GPS access while the app is in the foreground might be considered low‑risk, whereas the same access in the background would be high‑risk. This contextual weighting can partially compensate for missing log entries. We implemented a basic version of this fusion in a branch, but it was not included in the main evaluation because it complicated the baseline logic. In future releases, we plan to fully integrate it and test its effectiveness on fragmented devices.

Third, the version adaptation layer is a software engineering solution that abstracts away the API differences between Android versions. For instance, Android 12 introduced `getHistoricalOps(uid, packageName, attributionTag, ops, flags)` with an additional `attributionTag` parameter, while Android 13 added a `filter` flag to exclude system apps. Our adaptation layer detects the runtime version and invokes the appropriate method, ensuring that the data reading logic remains uniform. This layer is thoroughly unit‑tested with mocked API responses for each version, but we acknowledge that mocking cannot capture all vendor‑specific divergences. To further improve robustness, we employ a defensive programming strategy: if the API call returns an empty or malformed result, we retry with a lower `flags` value (e.g., `HISTORICAL\_OP\_FLAGS\_SINCE\_BOOT`). If all attempts fail, we log a warning and proceed with the last known good data, issuing a notification to the user that the audit may be incomplete.

Beyond the three strategies listed, we also consider user education as a mitigation. Many fragmentation issues are not technical but perceptual; users may not realise that their device's vendor has altered Android's default behaviour. By providing clear, non‑technical explanations in the app's onboarding flow (e.g., ``Your device manufacturer may limit permission history. To improve accuracy, please allow the app to run in the background.''), we set realistic expectations and encourage users to take the necessary steps. This human‑centred approach complements our technical mitigations and is consistent with the transparency and trust principles discussed in Section 6.3.

Finally, we must acknowledge that fragmentation is an ongoing challenge in the Android ecosystem, and no single framework can fully overcome it. Our work is a step towards addressing the problem, but it is not a panacea. We recommend that any production‑grade deployment of our framework include a telemetry module that anonymously reports the effectiveness of each mitigation strategy across different device models. Such data could be used to refine the adaptation layer and to alert users when their device is known to have severe limitations. This would turn fragmentation from a threat into an opportunity for continuous improvement, aligning with the iterative nature of privacy regulation itself.

\subsection*{Concluding Remarks on Limitations and Future Research}

In aggregate, the limitations discussed above point to a clear research agenda. First, the verification‑friction trade‑off could be explored through longitudinal user studies that measure actual privacy fatigue and behavioural responses to different notification frequencies. Second, the modular architecture could be extended with a formal verification framework that statically checks the consistency of baseline updates and module interactions. Third, the random testing methodology could be augmented with property‑based testing (using frameworks like KotlinTest) to systematically explore the input space. Fourth, the physical reality gap calls for a dedicated cross‑device testing lab, perhaps using a cloud‑based device farm, to collect empirical data on API behaviour and log retention policies across hundreds of Android models. Fifth, the fragmentation challenges could be mitigated by collaborating with device manufacturers to standardise permission log retention, though this is a long‑term policy goal beyond the scope of a single thesis.

Despite these limitations, our framework represents a solid foundation for user‑centric, regulation‑aware privacy supervision. Its emphasis on transparency, low cognitive friction, and statistically grounded decision‑making offers a viable alternative to both overly intrusive real‑time monitors and passive, non‑actionable permission dashboards. We believe that the integration of a rigorous random testing harness sets a new standard for validation in privacy tools, and we encourage other researchers to adopt similar methodologies to build trust in their systems.
